% =============================================================
% §1 引言
% =============================================================
\section{引言}
\label{sec:intro}

以大语言模型（LLM）推理为主导的生成式 AI 负载---亦称 AI 生成内容（AIGC）负载---已成为现代数据中心增长最快的算力类目：其能耗成本在五年间增长了数个数量级~\cite{patterson2021carbon}，且推理消耗的累计 FLOPs 已超过训练~\cite{patel2024splitwise}。随着 LLM 服务进入延迟敏感应用，系统社区正在重新思考推理负载应如何在异构算力资源池上\emph{调度}。

本文研究\textbf{云--边协同 LLM 推理调度}。该部署范式因三点原因日益受到关注：（i）\textbf{成本}---常见请求由更便宜的边缘加速器服务，云端 GPU 留给尾部请求；（ii）\textbf{隐私}---尽可能将 prompt 数据保留在本地；（iii）\textbf{尾延迟}---借助边缘临近性摊薄首 token 延迟。这里的\emph{协同}指部署层面的协作---云与边两层共同服务同一负载---而非多智能体决策：本文研究的每个调度器均为单一集中式策略。我们的优化目标（\S\ref{sec:sched-problem}）直接刻画尾延迟与成本：前者经由服务等级目标（\slo{}）达成率，后者经由每 token 能耗---最主要的边际服务成本；隐私是\emph{部署}云边资源池的动机，但与逐请求放置正交（可编码为可行性掩码，\S\ref{sec:mdp}，留作未来工作）。调度问题---决定每个请求应由哪台服务器服务---由此成为关键性能杠杆，且区别于 vLLM~\cite{kwon2023vllm}、Orca~\cite{yu2022orca}、Sarathi-Serve~\cite{agrawal2024sarathi} 等服务器内部批处理引擎。

\subsection{LLM 推理调度的特殊性}

现有集群调度器---异构最早完成时间（Heterogeneous Earliest Finish Time, HEFT）等表调度启发式、遗传算法（GA）与粒子群优化（PSO）等元启发式，以及 Decima~\cite{mao2019decima}、RLScheduler~\cite{zhang2020rlscheduler} 等学习式变体---原本为 Spark 作业、高性能计算（HPC）批任务等通用有向无环图（DAG）负载设计。它们并未建模从根本上区分 LLM 推理的五项物理特性：

\begin{enumerate}
\item \textbf{模型权重驻留。}LLM 权重（LLaMA-7B/13B/70B~\cite{touvron2023llama} 为 14--70~GB）持久驻留于 GPU 显存；切换模型需付出 5--60~s 的冷加载~\cite{fu2024serverlessllm}---比任何单请求计算高出数个数量级。

\item \textbf{KV cache 本地性。}每个请求的键值（KV）缓存与运行 prefill 的 GPU 绑定；迁移它需为每个请求移动数百 MB 数据，时常主导 decode 延迟~\cite{patel2024splitwise,kwon2023vllm}。

\item \textbf{连续批处理。}迭代级批处理以 ${\approx}(1 + (N-1)\alpha)\,T_{\text{solo}}$ 的时间服务 $N$ 个同模型请求---带来 $5$--$8\times$ 吞吐增益~\cite{yu2022orca,kwon2023vllm}，但只有当调度器将同模型请求共置时才能兑现。

\item \textbf{显存带宽受限的 decode。}与计算受限的 prefill 不同，decode 受高带宽显存（HBM）带宽约束：A100 级 GPU 上每 token ${\approx}20$~ms~\cite{patel2024splitwise,agrawal2024sarathi}，与算力无关。

\item \textbf{两阶段请求生命周期。}每个请求分解为 Prefill$\,\to\,$Decode，二者具有严格依赖与共享 KV 状态~\cite{zhong2024distserve}，违反了大多数 DAG 调度器的"各阶段独立"假设。
\end{enumerate}

这五项物理特性带来了通用负载所没有的权衡：集中同模型请求（利于批处理）还是分散它们（避免争用）；将 decode 派往快但远的服务器，还是派往其兄弟 prefill 服务器（免去 KV 迁移）。\textbf{忽略 AIGC 物理的调度器无法利用这些权衡}。

\subsection{为什么朴素方法会失效}

一个自然的猜想是：现成的强化学习（RL）调度器可以从基础特征（算力、队列长度、显存）中\emph{隐式}学到 AIGC 感知行为。我们正是为检验这一猜想而开展本工作；实证图景更为微妙（图~\ref{fig:pareto}）：

\begin{itemize}
\item 通用负载均衡基线（LeastLoaded、ShortestQueue）、最优 DAG 启发式（HEFT、GA、PSO）、PerLLM 风格的受限 bandit 调度器~\cite{yang2024perllm}，乃至现代深度强化学习（DRL）基线（A3C-R2N2~\cite{tuli2022a3cr2n} 与我们自实现的 Decima 风格图神经网络（GNN）变体~\cite{mao2019decima}）全部聚集在同一 Pareto 区域：\slo{} 达成率约 19--24\%，能效约 2.8--3.0 J/token。

\item 我们的 AIGC 感知近端策略优化（PPO）调度器与 A3C-R2N2 共用 actor-critic 家族与等价预训练预算，但加入显式 AIGC 状态特征与奖励整形，取得\textbf{明显更优的 Pareto 点}：\slo{} 30.2\%（相对 $+28\%$）、能耗 2.61~J/token（相对 $-7\%$），\slo{} 维度对全部 5 个最强单目标基线均 $p<0.05$。总体最强基线---我们同样实现的膝点 NSGA-II 搜索---在典型工况点接近本文的 \slo{}，但在中等负载下被显著超越（$p<0.01$），在每一种配置下付出更多能耗，且派发时计算开销高 $5\times$。

\item 然而，我们的 12 组件消融（表~\ref{tab:ablation}、图~\ref{fig:ablation}）揭示了一个悖论：移除任何单一 AIGC 组件（warm 奖励、batch 奖励、affinity 奖励、AIGC 状态特征，甚至全部联合移除）后，性能几乎不变。只有连续批处理仿真与充足预训练凸显为显著组件。
\end{itemize}

\subsection{调和与贡献}

认识到\textbf{AIGC 感知调度是一项系统级整体贡献，而非某个巧妙的奖励组件}，上述观察即可得到调和。一旦调度器具备（i）环境中的连续批处理物理、（ii）充足的预训练、（iii）状态或奖励中某种形式的 AIGC 感知，PPO 主干便能从任一 AIGC 先验中提取等价信号---它们可互相替代。Pareto 改进是真实的，但它来自\emph{组合}，而非单个奖励技巧。

\smallskip
\noindent 本文做出四项贡献：

\noindent\textbf{（C1）AIGC 推理仿真平台}，建模全部五项物理特性：带最近最少使用（LRU）驱逐的权重驻留、带 KV cache 显存核算的 prefill/decode 两阶段任务、准入时连续批处理，以及按 vLLM 实测标定的显存带宽下限。平台随复现 manifest 一并开源。

\noindent\textbf{（C2）AIGC 感知 PPO 调度器（RL）}，通过状态特征（已加载模型指示、批占用率、兄弟服务器亲和、KV cache 体积）与奖励整形（warm、batch、affinity bonus；cloud-overuse 惩罚）将 AIGC 物理暴露给策略网络，在 \slo{} 达成率与能效两个维度对 10 个强基线取得\textbf{均值性能上的 Pareto 占优}。

\noindent\textbf{（C3）系统性的 12 组件消融研究}，揭示哪些 AIGC 抽象真正有贡献。醒目的零结果---没有任何单项 AIGC 奖励或状态组件个体显著---警示后续工作：在缺乏联合消融时，不应将增益归因于特定奖励设计。

\noindent\textbf{（C4）实证的"工况"刻画}，表明 AIGC 感知 RL 仅在\textbf{中等负载工况}（$\lambda \in [1, 2]$ 请求/秒、uniform 模型混合）取得严格 Pareto 占优，同时在全部 30 种测试配置下保持最低的每 token 能耗，且相对最强基线的优势随负载增长---这一细致发现与 DRL 调度文献中常见的"处处取胜"论断相悖。

\smallskip
本文其余部分依次介绍背景（\S\ref{sec:background}）、形式化问题（\S\ref{sec:formulation}）、仿真器与调度器设计（\S\ref{sec:system}）、评估（\S\ref{sec:eval}）、讨论（\S\ref{sec:discuss}）、相关工作（\S\ref{sec:related}）与结论（\S\ref{sec:conclusion}）。
